% =====================================================================
%  Background: Iris and the logic of bunched implications (BI)
%  \input-able frame(s) only; no preamble. See main.tex for macros.
% =====================================================================

\begin{frame}{Iris, and the logic of bunched implications (BI)}
  The separation logic is built on the Lean port of \alert{Iris}; its
  assertions \(\mathrm{Prop}\) form a \wmech{BI (bunched-implication) logic}.
  \vfill
  % --- The assertion-language signature -------------------------------
  \begin{center}
  \renewcommand{\arraystretch}{1.35}
  \(
    \begin{array}{r@{\;\;}c@{\;\;}l@{\qquad}l}
      \ulcorner\,\cdot\,\urcorner
        & : & \mathrm{Prop} \to \mathrm{Prop}
        & \text{\footnotesize pure embedding} \\
      \mathsf{True},\ \mathsf{False},\ \mathsf{emp}
        & : & \mathrm{Prop}
        & \text{\footnotesize units} \\
      \wedge,\ \vee,\ \Rightarrow,\ \wprob{\ensuremath{\ast}},\ \wprob{\ensuremath{\wand}}
        & : & \mathrm{Prop} \times \mathrm{Prop} \to \mathrm{Prop}
        & \text{\footnotesize binary connectives} \\
      \forall,\ \exists
        & : & \forall A.\ (A \to \mathrm{Prop}) \to \mathrm{Prop}
        & \text{\footnotesize quantifiers}
    \end{array}
  \)
  \end{center}
  \vfill
  \begin{flushright}\footnotesize
    \wprob{\(\ast\)} separating conjunction \quad
    \wprob{\(\wand\)} magic wand \quad
    --- the separation-logic-specific connectives.
  \end{flushright}
\end{frame}


\begin{frame}{Iris, and the logic of bunched implications (BI)}
  Assertions \(\mathrm{Prop}\) form a \wmech{BI (bunched-implication) logic} on
  the Lean port of \alert{Iris}.
  \smallskip
  % --- The assertion-language signature -------------------------------
  \begin{center}
  \renewcommand{\arraystretch}{1.2}
  \(
    \begin{array}{r@{\;\;}c@{\;\;}l@{\qquad}l}
      \ulcorner\,\cdot\,\urcorner
        & : & \mathrm{Prop} \to \mathrm{Prop}
        & \text{\footnotesize pure embedding} \\
      \mathsf{True},\ \mathsf{False},\ \mathsf{emp}
        & : & \mathrm{Prop}
        & \text{\footnotesize units} \\
      \wedge,\ \vee,\ \Rightarrow,\ \wprob{\ensuremath{\ast}},\ \wprob{\ensuremath{\wand}}
        & : & \mathrm{Prop} \times \mathrm{Prop} \to \mathrm{Prop}
        & \text{\footnotesize binary connectives} \\
      \forall,\ \exists
        & : & \forall A.\ (A \to \mathrm{Prop}) \to \mathrm{Prop}
        & \text{\footnotesize quantifiers}
    \end{array}
  \)
  \end{center}
  \smallskip
  \begin{itemize}
    \item<2->
      \alert{A logic inside a logic.} \(\mathrm{Prop}\) is an object logic
      inside Lean's meta-logic.\footnote<3->{Revisited later in the talk.}
    \item<3->
      \alert{Impredicativity.} \(\forall/\exists\) range over an
      \emph{arbitrary} type \(A\) --- including \(\mathrm{Prop}\) itself: we
      may quantify over propositions.
  \end{itemize}
\end{frame}
